% =====================================================================
%  §8 Homological mirror conjecture
%  （原文 p52 第 18 行起，至 p58 末；§9 自 p59 起）
% =====================================================================
\section{同调镜像猜想}

在上一节中，我们构造了一个 $\Ainf$-范畴 $C(Y)$，它 $\Ainf$-等价于域 $\C_{\eps}$ 上
某个 Calabi--Yau 流形的凝聚层导出范畴。本节要构造一串 $\Ainf$-预范畴与
$\Ainf$-等价（参看 §6.6）
\[
  C_{unram}(Y)\to C_{unram,0}(Y)\leftarrow C^{tr}_{unram,0}(Y)
  \leftarrow C^{tr,\Pi}_{unram,0}(Y)\leftarrow FO(X^{\vee}),
\]
以及一个函子 $F:C_{unram}(Y)\to C(Y)$，它给出 $C_{unram}(Y)$ 与 $C(Y)$ 的某个满子
范畴之间的等价。回忆一下，本文所定义的 Fukaya--Oh 范畴 $FO(X^{\vee})$ 也等价于
Fukaya 范畴 $F(X^{\vee})$ 的一个满子范畴。于是我们便在 Fukaya 范畴 $F(X^{\vee})$
的满子范畴与 $C(Y)$ 的满子范畴之间建立了一个 $\Ainf$-等价。

我们要用的进路，与 §6 中处理 Morse 理论时所用的进路完全类似。上述
$\Ainf$-等价函子链中的所有范畴都具有同一类对象，即与 Fukaya--Oh 范畴相同的对象类。

\begin{yizhu}[链中的上标 $tr$ 与 $\Pi$]
这条链与 §6.6 中由 $M(Y)$、$DR^{tr,\Pi}_0(Y)$、$DR^{tr}_0(Y)$、$DR_0(Y)$、$DR(Y)$
构成的链逐项对应。其中 $tr$ 表示把对象序列限制在横截性受控的那些序列上所得的满子
范畴，$\Pi$ 表示再施行一次投影算子（Karoubi 型）完备化，即用投影算子
$\Pi_{f_0,f_1}$ 的像取代原来的态射空间；两种操作都不改变对象类，只改变态射空间。
原文对这两个上标未作说明，读者可比照 §6.6。
\end{yizhu}

\subsection{对象上的镜像对称函子}

这里我们将定义 dg-范畴 $C_{unram}(Y)$，以及它到 $C(Y)$ 的满忠实嵌入 $F$。

在附录中我们将解释复情形下镜像对称函子的通常图景。那里我们要用到一种沿环面纤维化的
纤维所作的 Fourier--Mukai 变换，这一变换的核是 Poincaré 丛的类比物。如果从一个
局部系出发，该局部系位于 $p^{\vee}:X^{\vee}\to Y$ 的某个拉格朗日截面上，那么此变换
把它变成一个 $X$ 上的光滑丛，其联络在反全纯方向上平坦。换言之，得到 $X$ 上的一个
全纯丛。

这些考虑在非阿基米德情形下不能逐字照搬。我们将按如下方式构造函子 $F$。设
$(L,\rho)$ 是范畴 $FO(X^{\vee})$ 的一个对象，满足 $\rank(\rho)=1$，且投影
$L\to Y$ 是一一映射。流形 $L$ 局部地由 $df$（模 $(\TYZ)^{\vee}$）的图像给出，其中
$f$ 是 $Y$ 上的光滑函数。对这样的对象，我们配置一个秩一的 $\OO_Y$-模层
$F(L,\rho)$。对充分小的开集 $U\subset Y$ 与选定的 $f\in C^{\infty}(U)$，层
$F(L,\rho)|_U$ 与 $\OO_Y|_U$ 等同。把 $f$ 换成 $f+l$（其中 $l\in Af_Y(U)$）会使
$F(L,\rho)|_U$ 的平凡化按 $\mathbf{1}_U\mapsto\exp(l)\mathbf{1}_U$ 改变（这里
$\mathbf{1}_U\in\OO_Y(U)$ 是恒等函数）。若 $\rank(\rho)$ 大于一，则对小的
$U\subset Y$ 把 $\rho|_U$ 分解为秩一局部系之和，再施行上述构造。类似地，若覆盖
$L\to Y$ 有多于一个叶，则对覆盖的每个叶施行上述构造，然后取直和。

\begin{yizhu}[$Af_Y$ 与平凡化的改变]
$Af_Y$ 是 $Y$ 上的局部仿射函数层：$Af_Y(U)$ 由形如 $\sum_i m_iy_i+c$ 的函数组成，
其中 $m_i\in\Z$、$c\in\R$，$y_i$ 是整仿射坐标（见 §5.2）。因此 $f\mapsto f+l$
正是把拉格朗日截面的生成函数作整仿射平移；它不改变 $L$，只改变层的平凡化。由于
$l$ 的斜率是整数，$\exp(l)$ 是 $\OO_Y$ 中的可逆元，故不同的 $f$ 给出同构的层，
其差别正是这一因子。
\end{yizhu}

我们把 $F$ 称为对象上的镜像对称函子。范畴 $C_{unram}(Y)$ 定义为这样的 dg-范畴：
其对象类为 $\Ob(FO(X^{\vee}))$，态射空间为
\[
  \Hom_{C_{unram}(Y)}\bigl((L_1,\rho_1),(L_2,\rho_2)\bigr)
  :=\Hom_{C(Y)}\bigl(F(L_1,\rho_1),F(L_2,\rho_2)\bigr).
\]

\subsection{态射的谱与半群}

设 $E_i=F(L_i,\rho_i)$，$i=1,2$，是对应于对象 $(L_i,\rho_i)\in FO(X^{\vee})$
（$i=1,2$）的局部自由 $\OO_Y$-模（即向量丛）。对任意
$\alpha\in\Hom_{C(Y)}(E_1,E_2)$ 与一点 $y\in Y$，我们把 $\alpha$ 在 $y$ 处的
\emph{谱}（spectrum）定义为某个（至多可数的）实数的离散集，带有限重数。

先设 $\rho_i$（$i=1,2$）是 $L_i$ 上的平凡秩一局部系，且 $L_i$（$i=1,2$）是 $Y$ 的
非分支（unramified）覆盖。对含 $y$ 的充分小开集 $U$，在局部坐标下可以写
$L_i=\mathrm{graph}(df_i)$（模 $(\TYZ)^{\vee}$），$i=1,2$，其中 $f_i:Y\to\R$ 是
光滑函数。态射 $\alpha\in\Hom_{C(Y)}(E_1,E_2)(U)=\widehat{\Omega}^{*}_Y(U)$ 在
$U$ 上的限制可以等同于无穷级数
\[
  \alpha=\sum_{i_1,\dots,i_n}c_{i_1\dots i_n}z_1^{i_1}\cdots z_n^{i_n},
  \qquad\text{其中}\quad
  c_{i_1\dots i_n}=\sum_j c_{j,i_1\dots i_n}e^{-\lam_{j,i_1\dots i_n}/\eps},
\]
而 $c_{j,i_1\dots i_n}\in\Omega^{*}_Y(U)$。

我们把 $\alpha$ 在 $y\in U$ 处的谱定义为实数集（带重数）
\[
  Sp_y(\alpha)
  =\Bigl\{-\lam_{j,i_1\dots i_n}+\sum_{1\le k\le n}i_ky_k+f_1(y)-f_2(y)\Bigr\},
\]
其中 $c_{j,i_1\dots i_n}$ 在 $y$ 处的芽不为零。可以验证 $Sp_y(\alpha)$ 是良定义的
（即不依赖于局部平凡化的选取），并且其唯一的极限点为 $s=-\infty$。

\begin{yizhu}[谱的读法]
谱把 $\alpha$ 的每一项按「指数」$-\lam_{j,i_1\dots i_n}+\sum_k i_ky_k$ 排成实数，
再统一平移 $f_1(y)-f_2(y)$；由于级数中 $\lam_j\to+\infty$，这些实数只在 $-\infty$
处堆积，故谱是离散的、重数有限的。这与辛侧按作用量给态射分层的做法互为类比：紧接着
的注给出的 $\R$-滤过只是同一信息的序化形式，两者互相决定。
\end{yizhu}

在局部系秩更高、且拉格朗日流形是 $Y$ 的非分支覆盖的一般情形，我们在 $y\in Y$ 附近
把 $E_i$（$i=1,2$）局部地分解为平凡秩一 $\OO_Y$-模的直和。态射在点 $y$ 处的谱于是
定义为相应线丛之间态射的谱的并。

\begin{rmk}
可以不用谱，而在态射空间上使用一个 $\R$-滤过
$\Hom_{C(Y)}(E_1,E_2)^{\le s}$。它来自态射层
$\underline{\Hom}_{\OO_Y}(E_1,E_2)\widehat{\otimes}\widehat{\Omega}^{*}_Y$
（完备张量积）的茎上的滤过，后者由条件
$\{-\lam_{j,i_1\dots i_n}+\sum_{1\le k\le n}i_ky_k+f_1(y)-f_2(y)\}\le s$ 定义。
容易看出，$\alpha$ 属于 $\Hom_{C(Y)}(E_1,E_2)^{\le s}$ 当且仅当对所有 $y\in Y$
都有 $Sp_y(\alpha)\subset(-\infty,s]$。
\end{rmk}

考虑代数态射构成的子空间
$\Hom^{alg}_{C(Y)}(E_1,E_2)\subset\Hom_{C(Y)}(E_1,E_2)$。它由有限和构成（对 $z_i$
与 $e^{-\lam_{j,i_1\dots i_n}/\eps}$ 都是有限和）。它在全体态射的空间中稠密（解析
函数可用 Laurent 多项式逼近）。此外，$\Hom_{C(Y)}(E_1,E_2)$ 与
$\Hom^{alg}_{C(Y)}(E_1,E_2)$ 关于上述 $\R$-滤过的完备化重合。

存在单参数半群（semigroup）$\phi^t$（$t\ge0$）作用在
$\Hom^{alg}_{C(Y)}(E_1,E_2)$ 上。在局部坐标下，$\phi^t$ 沿 $f_1-f_2$ 的梯度流
移动系数 $c_{j,i_1\dots i_n}$。为了整体地定义它，我们需要用几何的语言描述空间
$\Hom^{alg}_{C(Y)}(E_1,E_2)$。这将在下面完成。

给定上述两个拉格朗日子流形 $L_i\subset X^{\vee}$（$i=1,2$）、一点 $y\in Y$，以及
两点 $x_i\in L_i$（$i=1,2$）使 $p^{\vee}(x_i)=y$，我们定义一个集合
$P(L_1,L_2,y)$，其元素是路径 $\gamma\in(p^{\vee})^{-1}(y)$ 的同伦类，这些路径从
$x_1$ 出发、到 $x_2$ 终止。

每个同伦类含有环面上平坦度量下唯一的一条测地线。我们定义空间
$P(L_1,L_2)=\bigsqcup_{y\in Y}P(L_1,L_2,y)$。它带有显然的拓扑，使自然投影
$\pi:P(L_1,L_2)\to Y$ 成为具有可数纤维的非分支覆盖。利用 $X^{\vee}$ 上的辛形式
$\omega$，由公式 $\mu=\int_{\gamma}\omega$ 定义 $P(L_1,L_2)$ 上的闭 $1$-形式
$\mu$。在 $Y$ 上局部地有 $L_i=df_i$（模 $(\TYZ)^{\vee}$），$i=1,2$，其中
$f_i:Y\to\R$ 是光滑函数。于是在 $P(L_1,L_2)$ 上局部地有
$\mu=d(f_1-f_2+l)$，其中 $l$ 是层 $Af_Y$ 的拉回的局部截面。显然，函数 $l$ 在相差
一个实常数的意义下确定，于是得到 $P(L_1,L_2)$ 上的一个 $\R$-挠子
（$\R$-torsor）。利用嵌入 $\R\to\C^{*}_{\eps}$，$\lam\mapsto\exp(\lam/\eps)$，
得到 $\C^{*}_{\eps}$-挠子，它定义了 $P(L_1,L_2)$ 上的一维 $\C_{\eps}$-模构成的
局部系 $\C^{tw}_{\eps}$。$\C^{tw}_{\eps}$ 的纤维带有自然的滤过。事实上，在点
$(x_1,x_2,\gamma,y)\in P(L_1,L_2)$ 的一个邻域内可以选取光滑函数
$f=f_1-f_2+l$ 使 $\mu=df$，它给出 $\C^{tw}_{\eps}$ 的一个局部平凡化。在此平凡化
下，滤过对 $h\in\C_{\eps}$ 由条件 $v(h)(y)+f(y)\le s$（$s\in\R$）定义，其中 $v$
是赋值。我们定义 $\C^{tw}_{\eps}$ 的子层 $\C^{tw,alg}_{\eps}$，要求它在局部平凡化
下是「指数项的有限和」构成的子层。

注意有自然投影 $pr_i:P(L_1,L_2)\to L_i$，$i=1,2$。有了 $L_i$ 上的局部系 $\rho_i$
（$i=1,2$），我们便用关于 $pr_i$（$i=1,2$）的拉回，定义 $P(L_1,L_2)$ 上的局部系
$\widehat{\rho}_i$（$i=1,2$）。

在 $P(L_1,L_2)$ 上我们如下定义层 $\underline{\Hom}^{alg}(E_1,E_2)$（$E_i$，
$i=1,2$ 已在前面定义）：
\[
  \underline{\Hom}^{alg}(E_1,E_2)
  =\C^{tw,alg}_{\eps}\otimes(\widehat{\rho}_1)^{*}\otimes\widehat{\rho}_2
  \otimes\underline{\Omega}^{*}_{P(L_1,L_2)},
\]
其中 $\underline{\Omega}^{*}_{P(L_1,L_2)}$ 是微分形式层。我们给
$\underline{\Hom}^{alg}(E_1,E_2)$ 的茎配上 $\R$-滤过：它由 $\C^{tw}_{\eps}$ 上的
滤过诱导，其余张量因子取平凡滤过。

设 $\pi_!$ 表示紧支集直接像函子。则
\[
  \pi_!\bigl(\underline{\Hom}^{alg}(E_1,E_2)\bigr)
  =\pi_!\bigl(\C^{tw}_{\eps}\otimes\widehat{\rho}^{*}_1\otimes\widehat{\rho}_2
  \bigr)\otimes\underline{\Omega}^{*}_Y,
\]
其中最后一个张量因子是 $Y$ 上的 de Rham 微分形式层。

我们可以把 $\Z^n$ 与 $H_1(T^n,\Z)$ 等同，后者自然作用在路径 $\gamma$ 的同伦类上。
另一方面，$\Z^n$ 在 $\C_{\eps}$ 上的群环可以等同于 Laurent 多项式环
$\C_{\eps}[z_1^{\pm1},\dots,z_n^{\pm1}]$。设 $\C^{alg}_{\eps}\subset\C_{\eps}$ 是
指数项有限和构成的子环。容易看出，$\underline{\Hom}^{alg}(E_1,E_2)$ 的截面上
$\C^{alg}_{\eps}[\Z^n]$-模的结构，对应于它在 $\pi_!$ 下的像上
$\C^{alg}_{\eps}[z_1^{\pm1},\dots,z_n^{\pm1}]$-模的结构。利用这一观察可以证明
\[
  \Hom^{alg}_{C(Y)}(E_1,E_2)\simeq
  \Gamma\bigl(Y,\pi_!(\underline{\Hom}^{alg}(E_1,E_2))\bigr)
  =\Gamma_c\bigl(P(L_1,L_2),\underline{\Hom}^{alg}(E_1,E_2)\bigr),
\]
其中同构由层的自然态射
\[
  \pi_!\bigl(\underline{\Hom}^{alg}(E_1,E_2)\bigr)
  \to\underline{\Hom}^{alg}_{C(Y)}(E_1,E_2)
\]
诱导。这里 $\Gamma_c$ 指紧支集截面的函子。

利用 $Y$ 上的度量，我们给 $1$-形式 $\mu$ 配置 $P(L_1,L_2)$ 上的向量场 $\xi$。
局部地，$\xi$ 是 $f_1-f_2+l$ 的梯度流的生成元。不难证明，流的任何轨道都不会在
有限时间内跑向无穷远。因此向量场 $\xi$ 生成作用在 $P(L_1,L_2)$ 上的单参数半群
$\psi^t$。下述结果容易证明。

\begin{prp}
单参数半群 $\psi^t$ 在不属于 $L_1\cap L_2$ 的点的茎上使滤过下降。更确切地说，
\[
  \psi^t\bigl(\underline{\Hom}^{alg}_p(E_1,E_2)^s\bigr)
  \subset\underline{\Hom}^{alg}_{\psi^t(p)}(E_1,E_2)^{s-\int_0^t\mu},
\]
其中 $p\in P(L_1,L_2)$ 是任意一点。
\end{prp}

函子 $\pi_!$ 与层 $\underline{\Hom}^{alg}_{C(Y)}(E_1,E_2)$ 及
$\underline{\Hom}^{alg}(E_1,E_2)$ 的茎上的滤过相容。容易看出，前者之茎关于由
$\underline{\Hom}^{alg}(E_1,E_2)$ 上的滤过所诱导的滤过的完备化，与
$\Hom_{C(Y)}(E_1,E_2)$ 重合。由于半群 $\psi^t$ 使滤过下降，半群 $\phi^t$ 连续地
延拓到关于该滤过的完备化上。于是下述命题成立。

\begin{prp}
$\phi^t$ 的作用连续地从 $\Hom^{alg}_{C(Y)}(E_1,E_2)$ 延拓到
$\Hom_{C(Y)}(E_1,E_2)$。
\end{prp}

\subsection{阿贝尔簇的同调镜像对称}

进路与证明都与 §6 中的完全类似（只需把 §6 中的标量换成 $\C_{\eps}$）。细节从略。

在上一小节中，我们定义了作用在紧支集截面
$\Gamma_c(P(L_1,L_2),\underline{\Hom}^{alg}(E_1,E_2))$ 上的半群 $\psi^t$。这个
作用对应于半群 $\phi^t$ 在态射空间 $\Hom_{C(Y)}(E_1,E_2)$ 上的作用。注意，层
$\underline{\Hom}^{alg}(E_1,E_2)
=\C^{tw,alg}_{\eps}\otimes(\widehat{\rho}_1)^{*}\otimes\widehat{\rho}_2
\otimes\underline{\Omega}^{*}_{P(L_1,L_2)}$ 是层
$\C^{tw}_{\eps}\otimes(\widehat{\rho}^{*}_1\otimes\widehat{\rho}_2)
\otimes\underline{D}'_{P(L_1,L_2)}$ 的子层，其中
$\underline{D}'_{P(L_1,L_2)}$ 是 $P(L_1,L_2)$ 上的分布值微分形式层。后一层之
截面带有自然的拓扑：系数取分布的 $\exp(-\lam_i/\eps)$ 级数收敛，是指它与任一检验
微分形式配对后在 adic 意义下收敛。与 Morse 理论的情形（§6.5）类似，可以证明下述
结果。

\begin{prp}
对任意 $\beta\in\Gamma_c(P(L_1,L_2),\underline{\Hom}^{alg}(E_1,E_2))$，极限
\[
  \psi^{\infty}(\beta)=\lim_{t\to+\infty}\psi^t(\beta)\in
  \Gamma\bigl(P(L_1,L_2),\C^{tw}_{\eps}\otimes
  (\widehat{\rho}^{*}_1\otimes\widehat{\rho}_2)\otimes
  \underline{D}'_{P(L_1,L_2)}\bigr)
\]
存在。
\end{prp}

这一极限用生成 $\psi^t$ 的梯度流不难描述。利用 $\psi^t$ 把态射的谱移到
$-\infty$ 这一事实，与 §6.5 类似可以证明：极限 $\psi^{\infty}(\beta)$ 属于由相应于
不稳定流形 $U_x\subset P(L_1,L_2)$（$x\in L_1\cap L_2$）的分布所生成的有限维
$\C_{\eps}$-向量空间。显然，映射 $\beta\mapsto\psi^{\infty}(\beta)$ 延拓到关于
滤过的完备化上，并下降为映射 $\alpha\mapsto\phi^{\infty}(\alpha)$，其中
$\alpha\in\Hom_{C(Y)}(E_1,E_2)$。$\phi^{\infty}$ 的像落在同构于
$\Hom_{FO(X^{\vee})}((L_1,\rho_1),(L_2,\rho_2))$ 的空间中。

然后我们重复 §6.6 的论证。具体地说，我们按 §6.6 中定义范畴 $DR_0(Y)$ 的同样方式，
定义一个 $\Ainf$-预范畴 $C_{unram,0}(Y)$。$C_{unram,0}(Y)$ 的对象与
$C_{unram}(Y)$ 的相同。$C_{unram,0}(Y)$ 的态射空间是 dg-代数
$\C_{\eps}\widehat{\otimes}\Omega^{*0}$ 上的 dg-模，其中 $\Omega^{*0}$ 是微分形式
在辅助参数 $\delta\in\R_{\ge0}$ 的 $\delta=0$ 处的芽构成的 dg-代数（参看 §6.6）。
$C_{unram,0}(Y)$ 中态射的复合关于 dg-模结构是线性的。按定义，
$C_{unram,0}(Y)$ 与 $C_{unram}(Y)$ 中的横截性条件与 $FO(X^{\vee})$ 中的相同。于是
我们得到一个 $\Ainf$-预范畴 $C^{tr}_{unram,0}(Y)$，它 $\Ainf$-等价于
$C_{unram}(Y)$。

用同调微扰论，按 §6.6 的同样方式（此时投影算子与同伦都由半群 $\phi^t$ 定义），
我们构造范畴 $DR^{tr,\Pi}_0(Y)$ 的类比物。我们把这个 $\Ainf$-预范畴记为
$C^{tr,\Pi}_{unram,0}(Y)$。此范畴的态射空间是 $\Omega^{*0}$ 与有限维
$\C_{\eps}$-向量空间的完备张量积，后者由不稳定流 $[U_x]$ 的「光滑化」张成。这些
光滑化按 §6.5 的同样方式、借助算子 $R_{\delta}$ 定义。按定义，范畴
$C^{tr,\Pi}_{unram,0}(Y)$ 的横截性条件与范畴 $FO(X^{\vee})$ 的相同，其态射空间
自然拟同构于 $FO(X^{\vee})$ 中相应的态射空间。重复 §6.6 的论证会看到，
$C^{tr,\Pi}_{unram,0}(Y)$ 上的 $\Ainf$-结构与 $FO(X^{\vee})$ 上的
$\Ainf$-结构等价。事实上，从空间
$\Hom_{FO(X^{\vee})}((L_1,\rho_1),(L_2,\rho_2))$ 到空间
$\Hom_{C^{tr,\Pi}_{unram,0}(Y)}(F(L_1,\rho_1),F(L_2,\rho_2))$ 有自然映射。回忆
一下，范畴 $C_{unram}(Y)$、$C_{unram,0}(Y)$、$C^{tr}_{unram,0}(Y)$ 与
$C^{tr,\Pi}_{unram,0}(Y)$ 都具有同一类对象。因此我们有对象上的镜像对称函子
$F:FO(X^{\vee})\to C^{tr,\Pi}_{unram,0}(Y)$（见 §8.1）。另一方面，上述考虑给出
态射空间之间的线性映射。设 $\nu_{X_1,X_2}$ 是对两个对象 $X_i=(L_i,\rho_i)$
（$i=1,2$）的这种线性映射。下述命题的证明与 §6.6 中相应命题的证明完全类似。
（由于现在我们在域 $\C_{\eps}$ 上工作，要用 §6.2 最后一段中 $m_k$（$k\ge1$）的
公式。）

\begin{prp}
设 $E_i=F(X_i)$，$0\le i\le k$，$k\ge1$，是对应于对象
$X_i=(L_i,\rho_i)\in FO(X^{\vee})$（$0\le i\le k$）的秩一局部自由
$\OO_Y$-模（向量丛）。则对
\[
  m_k^{FO(X^{\vee})}:\otimes_{0\le i\le k}\Hom(E_i,E_{i+1})
  \to\Hom(E_0,E_k)[2-k]
\]
的公式，与对
\[
  m_k^{C^{tr,\Pi}_{unram,0}(Y)}:\otimes_{0\le i\le k}\Hom(X_i,X_{i+1})
  \to\Hom(X_0,X_k)[2-k]
\]
的公式，在把标量从 $\C_{\eps}$ 扩张到 $\C_{\eps}\widehat{\otimes}\Omega^{*0}$
之后，当态射空间通过映射 $\nu(X_i,X_j)$ 等同起来时，彼此吻合。
\end{prp}

\begin{yizhu}[原文的拼写]
原文在 p57 末与 p58 中把范畴 $C^{tr,\Pi}_{unram,0}(Y)$ 的其余几处写作
$C^{tr,\Pi}_{unrum,0}(Y)$（$unrum$ 应为 $unram$，即 unramified 之缩写），系排印
疏漏。译文统一作 $C^{tr,\Pi}_{unram,0}(Y)$，以免读者误以为这是另一个范畴。
\end{yizhu}

因此 $\Ainf$-预范畴 $C^{tr,\Pi}_{unram,0}(Y)$ 与 $FO(X^{\vee})$ 等价。另一方面，
由 §6.4、§6.6 可知 $C_{unram}(Y)$ 与 $C^{tr,\Pi}_{unram,0}(Y)$ 也等价。最后，
应用函子 $F$，我们得到下述定理。

\begin{thm}
$C(Y)$ 的满子范畴 $F(C_{unram}(Y))$ 与 $FO(X^{\vee})$ 是 $\Ainf$-等价的。
\end{thm}

这就是我们关于阿贝尔簇的同调镜像猜想的表述。

% 原文 p59 顶部（Remark 21，属 §8 收尾；§9 标题之前）
\begin{rmk}
若给环面 $Y=\R^n/\Z^n$ 配上平坦度量，并且只考虑 $X^{\vee}$ 中的平坦拉格朗日子环面，则
$\Ainf$-预范畴 $FO(X^{\vee})$ 中所有高次复合都可以用显式的「截断 theta 级数」
（truncated theta series）写出，与 \cite{Ko}、\cite{P1} 在椭圆曲线情形所考虑的那些级数类似。
\end{rmk}

\begin{yizhu}[定理 4 证明的性质]
§8.3 并未对定理 4 另给一套独立的证明。作者开宗明义地说「进路与证明都与 §6 完全
类似」、细节从略：实际的论证链是，命题 14 把 $FO(X^{\vee})$ 的 $\Ainf$-结构搬到
$C^{tr,\Pi}_{unram,0}(Y)$ 上，再由 §6.4、§6.6 的等价把后者与 $C_{unram}(Y)$ 连起来，
最后经函子 $F$ 得到结论。因此定理 4 是把 §6 的整套机器照搬过来、把标量换成
$\C_{\eps}$ 所得的推论；本节给出的是一条证明思路，其细节须回查 §6。
\end{yizhu}

% [原文 p52–p59 顶部]
